% Generated draft: requires mathematical and source review.
\documentclass[11pt]{article}
\usepackage[margin=25mm]{geometry}
\usepackage{fontspec}
\setmainfont{Latin Modern Roman}
\usepackage{amsmath,amssymb,mathtools,unicode-math}
\setmathfont{Latin Modern Math}
\usepackage{enumitem,needspace,xurl,hyperref}
\setlength{\parindent}{0pt}
\setlength{\parskip}{5pt}
\setlength{\emergencystretch}{3em}
\newcommand{\N}{\mathbb{N}}
\newcommand{\Z}{\mathbb{Z}}
\newcommand{\Q}{\mathbb{Q}}
\newcommand{\R}{\mathbb{R}}
\newcommand{\C}{\mathbb{C}}
\newcommand{\F}{\mathbb{F}}
\newcommand{\A}{\mathbb{A}}
\newcommand{\PP}{\mathbb{P}}
\newcommand{\E}{\mathbb{E}}
\newcommand{\eps}{\varepsilon}
\newcommand{\dd}{\,\mathrm{d}}
\begin{document}
\section*{UnsolvedMath ID: 30006425}
\textbf{Source problem code:} catalogue entry
\subsection{Definitions and mathematical statement}
Let \(F_2=\langle a,b\mid\ \rangle\) be the free group on two generators, and put \(D=C^*(F_2\times F_2)\). Here \(C^*(G)\) denotes the full group \(C^*\)-algebra of a discrete group \(G\), characterized by its universal property for unitary representations of \(G\).

For the unital algebra \(D\), use the following unital formulations of the lifting properties. A linear map is completely positive if all its matrix amplifications preserve positivity, and is unital completely positive if it also preserves the identity. Let \(A\) be a unital \(C^*\)-algebra. It has the \emph{lifting property} (LP) if, for every unital \(C^*\)-algebra \(B\), every proper closed two-sided ideal \(J\triangleleft B\), and every unital completely positive map \(\phi:A\to B/J\), there is a unital completely positive map \(\widetilde\phi:A\to B\) such that
\[
q\circ\widetilde\phi=\phi,
\]
where \(q:B\to B/J\) is the quotient map. It has the \emph{local lifting property} (LLP) if, for every such \(B,J,\phi\) and every finite-dimensional operator system \(E\subseteq A\), there is a unital completely positive map \(\widetilde\phi_E:E\to B\) such that
\[
q\circ\widetilde\phi_E=\phi|_E.
\]
An operator system here is a complex linear, self-adjoint subspace containing the identity, with the matrix order inherited from \(A\).

\textbf{Specific questions.} Determine separately whether
\[
D\text{ has the LLP},\qquad D\text{ has the LP}.
\]
The source notation ``(L)LP'' is retained as these two questions; the word ``or'' in the cleaned record is not interpreted as an exclusive alternative.

\textbf{Proposed general equivalence.} The record also asks whether
\[
A\text{ has the LLP}\quad\Longleftrightarrow\quad
\operatorname{Ext}(A)\text{ is a group}.
\]
The intended class of \(C^*\)-algebras \(A\) and the convention for \(\operatorname{Ext}(A)\) are not specified. A customary extension-semigroup setting involves extensions
\[
0\longrightarrow\mathcal K(H)\longrightarrow X\longrightarrow A\longrightarrow0,
\]
where \(H\) is a separable infinite-dimensional Hilbert space and \(\mathcal K(H)\) is its algebra of compact operators, with addition induced by direct sum. However, the record does not specify the equivalence relation, the treatment of split or essential extensions, or separability and unitality hypotheses. Under an extension-semigroup interpretation, being a group means that every extension class has an additive inverse relative to the neutral class. This description identifies the intended type of question but does not fix its missing conventions. In particular, the unital lifting definitions above do not silently restrict the general question to unital algebras.

\textbf{Supplied historical context and qualifications.} A \(C^*\)-algebra is residually finite-dimensional (RFD) if its finite-dimensional \(*\)-representations separate its elements. The original passage states that the Connes embedding problem is equivalent to \(D\) being RFD. The supplied dated assessment says that \(\mathrm{MIP}^*=\mathrm{RE}\) refuted the Connes embedding problem and consequently refuted this RFD assertion; it describes the LLP question and the proposed extension-semigroup equivalence as unresolved. These are supplied historical and literature assessments, not independent conclusions of this draft. The assessment does not separately resolve or explicitly assess the LP question.

The original passage additionally reports a theorem labeled [9]: if \(F\) is a finite-dimensional common \(C^*\)-subalgebra of \(A\) and \(B\), and both \(A\) and \(B\) have the LP, then their full amalgamated free product \(A*_F B\) has the LP. It contrasts this with an earlier statement of uncertainty attributed to [12, p.~15]. The bibliographic identities and full hypotheses of [9] and [12] are absent. The passage further claims consequences for finite tree products with finite edge groups and, in particular, finitely generated virtually free groups. No vertex-group hypotheses for the former claim are supplied, so it is not expanded into a universal mathematical assertion here. It also reports a new proof of Blackadar's semiprojectivity result for free products amalgamated over finite-dimensional subalgebras, a new proof of Li and Shen's RFD characterization for unital such free products, and a characterization in the non-unital case; the characterization criteria are not supplied. Its comment that few LP examples were then known beyond nuclear algebras and the stated constructions is historical source commentary, not a current classification.
\subsection{Short English statement}
The main questions ask separately whether the full group \(C^*\)-algebra \(C^*(F_2\times F_2)\) has local or global completely positive lifting properties. A further question proposes that LLP is equivalent to invertibility of all extension classes. The latter question needs a specified algebra class and extension-semigroup convention. The supplied partial-resolution label concerns a related RFD assertion and does not constitute a verified resolution of the lifting questions.
\subsection{Sources}
\begin{itemize}
\item Local registry: \texttt{data/unsolvedmath/problems.json}, with \texttt{catalog/problems.json} as the supplementary registry specified by the project; supplied complete source record, permanent UnsolvedMath ID \(30006425\), source code \texttt{catalogue entry}. No individual proposer is supplied. The record gives proposed year 2025.
\item Supplied report citation: \emph{C*-Algebras (2026)}, Oberwolfach Reports, DOI \texttt{10.4171/owr/2025/35}, pp.~1887--1948. \url{https://doi.org/10.4171/owr/2025/35}. The citation's title year 2026 differs from the proposed year 2025 and the year component of the DOI; these data are retained without reconciliation.
\item Supplied literature source: ``MIP-star equals RE.'' \url{https://arxiv.org/abs/2001.04383}.
\item Supplied literature source: ``Pop Ext and LLP results.'' \url{https://doi.org/10.7153/oam-2020-14-36}.
\item Supplied literature source: ``Oberwolfach Report 35.'' \url{https://doi.org/10.4171/owr/2025/35}.
\item The literature assessment was checked on \textbf{2026-08-22}, according to the supplied record. This is supplied dated curation, not independently verified. Its status label is \texttt{partially\_solved}; its assessment says that the related RFD assertion is false and that the LLP and general \(\operatorname{Ext}\) questions remain unresolved. The registry's \texttt{corrected\_verified} statement label is accompanied by a description limiting the repairs to delimiters, notation, grammar, and typography; it is not treated here as independent mathematical verification.
\item The original excerpt attributes contextual results to Blackadar and to Li and Shen, and uses reference labels [9] and [12, p.~15]. Their complete bibliographic entries are not supplied. No identification of these numbered references with the other supplied links is assumed.
\end{itemize}
\end{document}
